% Generated draft: requires mathematical and source review.
\documentclass[11pt]{article}
\usepackage[margin=25mm]{geometry}
\usepackage{fontspec}
\setmainfont{Latin Modern Roman}
\usepackage{amsmath,amssymb,mathtools,unicode-math}
\setmathfont{Latin Modern Math}
\usepackage{enumitem,needspace,xurl,hyperref}
\setlength{\parindent}{0pt}
\setlength{\parskip}{5pt}
\setlength{\emergencystretch}{3em}
\newcommand{\N}{\mathbb{N}}
\newcommand{\Z}{\mathbb{Z}}
\newcommand{\Q}{\mathbb{Q}}
\newcommand{\R}{\mathbb{R}}
\newcommand{\C}{\mathbb{C}}
\newcommand{\F}{\mathbb{F}}
\newcommand{\A}{\mathbb{A}}
\newcommand{\PP}{\mathbb{P}}
\newcommand{\E}{\mathbb{E}}
\newcommand{\eps}{\varepsilon}
\newcommand{\dd}{\,\mathrm{d}}
\begin{document}
\section*{UnsolvedMath ID: 30002978}
\textbf{Source problem code:} catalogue entry
\subsection{Definitions and mathematical statement}
Let \(\psi_{p,n,\delta}\) denote the universal \(K3\) moduli map named in the supplied source record, and let \(g\) denote the genus parameter appearing there. The record does not define the domain, target, admissible parameters, or the relation between \(g\) and \((p,n,\delta)\). Accordingly, write
\[
\psi_{p,n,\delta}\colon X_{p,n,\delta}\longrightarrow Y_{p,n,\delta}
\]
only as placeholder notation for the intended map; this does not specify either moduli space.

For a map \(f\colon X\to Y\), its image is \(\operatorname{Im}(f)=\{f(x):x\in X\}\). A divisorial image is usually understood to mean that its closure is a codimension-one locus in the target. Under that interpretation, the condition would be
\[
\operatorname{codim}_{Y_{p,n,\delta}}
\overline{\operatorname{Im}(\psi_{p,n,\delta})}=1.
\]
The source record does not specify whether closure is intended, whether an actual effective divisor is required, or how components are to be treated. Higher-rank Brill--Noether theory concerns loci of vector bundles of rank at least two satisfying conditions on their spaces of sections; the relevant ranks, degrees, stability conditions, and section bounds are not specified here.

\textbf{Question as supplied.} Among the admissible maps \(\psi_{p,n,\delta}\) with \(g\leq 11\), can any image be a divisor in its intended target? Can some images \(\operatorname{Im}(\psi_{p,n,\delta})\) be described by higher-rank Brill--Noether theory?

These are existence and description questions, not assertions that every such image is divisorial or has such a description. A fully specified mathematical formulation requires the missing moduli-map definitions and hypotheses from the cited report. The record also leaves unclear whether the restriction \(g\leq 11\) applies to the second question.
\subsection{Short English statement}
The problem asks whether universal \(K3\) moduli maps can have divisorial images for \(g\leq 11\), and whether some images admit higher-rank Brill--Noether descriptions. The supplied record omits the definitions and parameter restrictions needed to make these questions fully precise. Its partially solved label and genus-ten assessment are supplied dated curation, requiring independent mathematical review.
\subsection{Sources}
\begin{itemize}
\item Supplied local UnsolvedMath registry record, permanent ID \(30002978\), title \emph{Divisorial Images of Universal K3 Moduli Maps}, exact source problem code \texttt{catalogue entry}, in the Oberwolfach Reports Open Problems collection. No individual proposer is identified; the supplied proposed year is 2015.
\item Supplied report citation: \emph{Mini-Workshop: Singular Curves on \(K3\) Surfaces and Hyperkähler Manifolds} (2016), \emph{Oberwolfach Reports}, pp.\ 2939--2967, DOI \texttt{10.4171/owr/2015/51}. \url{https://doi.org/10.4171/owr/2015/51}. The record labels the statement \texttt{corrected\_verified} and reports checking it against the official report; that verification is supplied metadata, not an independent verification for this draft.
\item Farkas--Ortega, \emph{Higher-Rank Brill--Noether Theory}. Supplied literature reference: \url{https://arxiv.org/1102.0276}.
\item \emph{Moduli of Nodal Curves on K3 Surfaces}. Supplied literature reference: \url{https://arxiv.org/1502.07378}.
\item Knutsen--Torelli, \emph{Brill--Noether Loci on Severi Varieties}. Supplied literature reference: \url{https://arxiv.org/2411.14562}.
\item Literature assessment dated exactly \texttt{2026-08-21}: the supplied status is partially solved. The supplied assessment states that the genus-ten \(K3\) image has a higher-rank Brill--Noether realization and that related loci have been studied, while a general divisorial-image criterion for \(g\leq 11\) and systematic descriptions of all images remain open. This is supplied dated curation derived from the OWR record, not independently verified; the time-sensitive status and cited claims require independent verification.
\end{itemize}
\end{document}
